\chapter{関連研究}
\label{ch:related-work}

\section{逐次アルゴリズムの系譜}
\label{sec:related-sequential}

極大クリーク列挙の逐次アルゴリズムは、Bron--Kerbosch 法を起点とする一つの系譜をなす。
その技術的な内容は第\ref{ch:preliminaries}章で述べたとおりである。

Tomita, Tanaka, Takahashi は、ピボット選択によって重複した枝の展開を防ぐ CLIQUES を提案し、その最悪計算量が極大クリーク数の理論上限 $3^{n/3}$ と一致することを示した\cite{tomita2006}。
この上限自体は、Moon と Moser が示した極値グラフによって与えられている\cite{moon1965}。
Conte と Tomita は、CLIQUES と Bron--Kerbosch 法の各変種について、全体計算量だけでなく極大クリークを一つ出力するごとの遅延（delay）まで踏み込んで解析し直した\cite{conte2022}。
本論文の CLIQUES 実装は、この論文が与える擬似コードに従っている。

Eppstein, L\"offler, Strash は、CLIQUES のピボット選択自体は保ちながら、最外層だけを縮退順序に沿って処理することで、計算量を縮退数 $d$ を用いた $O(d n 3^{d/3})$ に抑える変種を示した\cite{eppstein2010}。
Eppstein と Strash は、続く実験で、実データを含む大規模な疎グラフでもこの変種が CLIQUES より高速に動作することを確かめている\cite{eppstein2013}。
本論文は、この Eppstein のアルゴリズムを、CLIQUES との逐次比較を経た上で並列化の対象に選んでいる。
選定の根拠となる測定結果は第\ref{ch:experiments}章で示す。

\section{並列化に関する先行研究}
\label{sec:related-parallel}

MCE の並列化についても複数の先行研究がある。

Schmidt, Hadfield, Samatova, Rice, Schuchardt は、共有メモリ環境で動作するスケーラブルな並列 MCE アルゴリズムを提案した\cite{schmidt2009}。
Das, Svendsen, Tirthapura による\textbf{ParMCE} は、共有メモリ上で動作する並列 MCE アルゴリズムであり、32 コア環境で逐次実行比最大 31 倍の高速化を報告している\cite{das2018parmce}。
グラフ分割に基づく別の手法も、既存手法比で最大 10 倍の高速化を報告している\cite{jpdc2017kplex}。
これらの論文はいずれも抄録レベルでしか内容を確認できておらず、具体的な実装は本文を未読のままである。
\todo{各論文の本文を確認し、アルゴリズムと実験条件の詳細を加筆する}

抄録から読み取れる範囲でも、これらの先行研究には共通する設計上の定石が見て取れる。
第一に、探索木をあらかじめプロセッサ数よりずっと多いサブツリーに分割しておく\textbf{オーバーデコンポジション}である。
第二に、暇になった worker が未処理のサブツリーを実行時に取りに行く\textbf{work stealing}（動的スケジューリング）であり、負荷の偏りを実行しながら解消する。
第三に、グラフ本体を worker ごとにコピー転送するのではなく共有メモリ上に置く配置であり、コピーにかかる固定費そのものを発生させない。

\section{本研究の位置づけ}
\label{sec:related-position}

本論文の並列実装は、これらの定石とは異なり、縮退順序で並べた頂点をバッチに区切って worker へ静的に配る素朴な配分から出発している\footnote{バッチの並べ方（block、reversed、interleave の3種）は第\ref{ch:implementation}章で述べる。}。
静的な配分は、work stealing のような実行時の負荷再分配を持たないため、負荷の偏りをバッチの作り方でしか吸収できない。
それでも、第\ref{ch:experiments}章で示すとおり、バッチの並べ方を変えるだけで worker の遊休はほぼ解消できた。
一方で、遊休を解消しても実行時間の短縮は先行研究が報告するような大幅な高速化には届かず、その理由は性能コアと効率コアの非対称な構成や共有資源の競合という、配分の工夫だけでは解決できない要因にあった。

先行研究は、これらの定石がなぜ有効かという機構までは抄録から読み取れない。
本論文は、単一マシン上での実測によって、work stealing とオーバーデコンポジションが解消しようとしている負荷の偏りの実体（重い頂点が縮退順序の末尾に密集すること）と、共有メモリ配置が解消しようとしている固定費の実体（プロセス起動時のグラフ転送コスト）を、それぞれ切り分けて確認する。
新しい並列アルゴリズムを提案するのではなく、既存の定石が効く理由とその限界を単一マシンの実測によって裏づける、再現・分析型の研究として本論文を位置づける。
